\scanpage{85}
\uline{Conjectures sur la cohomologie des schémas affines}

\margin{\unsure{1}}
$X$ : (``bon'') schéma affine intègre pt gén. $\unsure{\alpha|x|a}$.
$\ell$ \unsure{nb} premier

$1^\circ$) Si $\ell 1_X$ inversible, alors
\[
\operatorname{cd}_{\ell}(X)\leq\operatorname{cd}_{\ell}(k(x))
\]
\margin{\deleted{\unsure{[}} [et cela devient une égalité pour les ouverts assez petits ?]}

$2^\circ$) Si $\ell 1_X$ non inversible et non nilpotent, on a
\[
\operatorname{cd}_{\ell}(X)\leq\sup\bigl(\operatorname{cd}_{\ell}(k(x)),\unsure{\delta|\nu}(V(\ell))+3\bigr)
\]
où on pose, pour un schéma $Z$ de car. $\ell$
\[
\delta(Z)=\sup_{\substack{z\ \text{maximal}\\\text{de }Z}}\delta(k(z))
\]
et pour un corps $k$ de car $\ell$
\[
\delta(k)\quad\text{tel que}\quad \ell^{\delta(k)}=[k:k^\ell]
\]

\deleted{$1^\circ$ bis) Sous \unsure{la condition|les conditions} de $2^\circ$, lorsque $\exists\ \ell'$ \deleted{\unsure{un}} \inserted{premier} tel
que $\ell'1_X$ inversible i.e. $X$ sur $S_{\ell'}$ ($S=\operatorname{Spec}\mathbb Z$),
alors on peut remplacer \inserted{dans la formule}
$(\delta(V(\ell))+3)$ par $\delta(V(\ell))+1$. Sous les conditions
de $2^\circ$, lorsque}
% slanted
\margin{\deleted{\uline{Question} A-t-on tjrs $\operatorname{cd}_{\ell}(k(x))\unsure{\geq|>}\delta(V(\ell))+1$}}

Supposons $\exists$ \deleted{\unsure{un entier}}\inserted{$\ell'$}
premier tel que $\mathbb Z_{\ell'}$ inversible. Alors
\[
\operatorname{cd}_{\ell}(X)\deleted{\unsure{\leq\sup(}}\inserted{\leq}\operatorname{cd}_{\ell}(k(x))\deleted{\unsure{,\delta(V(\ell))+1)}}
\]

$4^\circ$) $A$ bon anneau local hensélien intègre, \unsure{cor.|car.} \unsure{résiduel|résiduelle} $\unsure{k|\ell}$,
car $K=0$. Soit $U$ ouvert affine de $X=\operatorname{Spec} A$. Alors
\[
\operatorname{cd}_{\ell}(U)\leq\dim A+\delta(k)+\operatorname{cd}_{\ell}(k)
\]
% archive sheet number: 84
